\section{Forelesning 5: \emph{Digital innovasjon, disrupsjon og strategier}}
\label{sec:lec05}

\subsection*{Oversikt}
Forelesning 5 (Sebastian Winslow) tar for seg \emph{digital innovasjon} som et
eget fenomen, distinkt fra alminnelig innovasjon, og kobler dette til
Christensens disrupsjonsteori. Hovedinnsikten: digitale komponenter har
egenskaper (homogenitet, redigerbarhet, selvreferering) som muliggjør stadig
\emph{rekombinasjon} -- og dermed nye former for innovasjon og disrupsjon.

\subsection{Hovedteori}

\subsubsection*{Definisjoner av innovasjon}
Pensum gir flere komplementære definisjoner:
\begin{itemize}
  \item \textbf{Rogers (1962/1983)}: «En idé, praksis eller gjenstand som
    oppfattes som ny av et individ eller en adopsjonsenheten.»
  \item \textbf{Hukal \& Henfridsson (2017)}: «Samskapelse av nye tilbud
    gjennom \emph{rekombinasjon} av digitale og/eller fysiske komponenter.»
  \item \textbf{Yoo et al.\ (2010)}: «Gjennomføring av nye kombinasjoner av
    digitale og fysiske komponenter for å produsere nye produkter.»
  \item \textbf{Nambisan et al.\ (2017)}: «Bruk av digital teknologi under
    innovasjonsprosessen.»
\end{itemize}

\begin{definition}{Digital innovasjon}
Gjennomføring av nye \emph{kombinasjoner} av digitale og fysiske komponenter
for å skape nye produkter, tjenester eller prosesser (Yoo et al.\ 2010).
\end{definition}

\subsubsection*{Tre egenskaper ved digital teknologi}
Det som gjør digital teknologi spesielt innovativ:
\begin{itemize}
  \item \term{Data homogeneity} (datahomogenitet) -- all digital informasjon
    er bits; ulike datatyper kan kombineres uten kompatibilitetsproblemer.
  \item \term{Editability} (redigerbarhet) -- digitalt innhold kan endres,
    også \emph{etter} det er distribuert.
  \item \term{Self-referentiality} (selvreferering) -- digitale systemer kan
    referere til og endre seg selv; programvare kan modifisere programvare.
\end{itemize}

\subsubsection*{Rekombinasjon}
Digital innovasjon = kombinere eksisterende digitale + fysiske komponenter på
nye måter. \keypoint{Innovasjon trenger ikke være «nytt fra grunnen av»; det
er ofte en \emph{ny kombinasjon} av eksisterende byggeklosser.}

\subsubsection*{Christensens disrupsjonsteori}

\begin{definition}{Disruptiv innovasjon}
En innovasjon som starter i bunnen av markedet eller i et forsømt segment
med et enklere, billigere tilbud som etablerte aktører ignorerer. Over tid
forbedrer disruptøren seg og beveger seg oppover, og fortrenger etablerte
aktører som har fokusert på sine beste kunder.
\end{definition}

Distinksjoner:
\begin{itemize}
  \item \term{Sustaining innovation} -- forbedrer eksisterende verdiforslag
    for eksisterende kunder.
  \item \term{Low-end disruption} -- enklere/billigere tilbud nedenfra.
  \item \term{New-market disruption} -- skaper et helt nytt marked av
    ikke-konsumenter.
\end{itemize}

\subsubsection*{Strategier mot/for digital disrupsjon}
Som i \autoref{sec:lec02}: forsvare, investere, kjøpe, samarbeide, selv-disrupsjon.
Christensens råd til etablerte: opprett en separat enhet med eget mandat for å
forfølge disruptive muligheter -- den etablerte organisasjonen vil naturlig
avvise dem.

\subsection{Sentrale fagbegreper}
\begin{itemize}
  \item \term{Digital innovation} -- rekombinasjon av digitale + fysiske
    komponenter.
  \item \term{Recombination} -- nye kombinasjoner av eksisterende byggeklosser.
  \item \term{Data homogeneity}, \term{editability}, \term{self-referentiality}
    -- de tre egenskapene.
  \item \term{Disruptive innovation} (Christensen).
  \item \term{Sustaining innovation} -- forbedrer eksisterende posisjon.
  \item \term{Low-end} og \term{new-market disruption}.
\end{itemize}

\subsection{Modeller og rammeverk}

\figureplaceholder{lec05-christensen-disruption-curve}{Christensens
disrupsjonskurve -- disruptøren starter i bunnen og beveger seg oppover.}

\figureplaceholder{lec05-digital-recombination}{Digital innovasjon som
rekombinasjon -- digitale + fysiske komponenter danner nye tilbud.}

\subsubsection*{Diagnostiske spørsmål for å vurdere et initiativ}
\begin{enumerate}
  \item Bruker initiativet eksisterende digitale + fysiske komponenter på en
    ny måte? $\rightarrow$ digital innovasjon.
  \item Forbedrer det det eksisterende verdiforslaget for eksisterende
    kunder? $\rightarrow$ sustaining.
  \item Treffer det en ikke-betjent eller forsømt kundegruppe med et
    enklere/billigere tilbud? $\rightarrow$ potensielt disruptivt.
\end{enumerate}

\subsection{Eksempler og caser}
\begin{itemize}
  \item \textbf{Netflix vs.\ Blockbuster} -- klassisk disrupsjon: Netflix
    startet med et enklere postordretilbud, beveget seg deretter oppover
    med streaming.
  \item \textbf{Coop App} -- \emph{sustaining innovation}, ikke disrupsjon.
    Den forsterker det eksisterende verdiforslaget (lojalitet,
    butikkopplevelse). Mulige disruptorer av Coop: rene
    netthandelsaktører som Nemlig; lavprisekspansjon (Lidl, Netto);
    plattformaktører som tilbyr matplanlegging/bestilling i stor skala.
  \item \textbf{Rekombinasjon i Coop App}: lojalitetsdata (digital) +
    butikkbesøk (fysisk) + personlige tilbud (digital) + Scan \& Betal
    (digital + fysisk) $\rightarrow$ et nytt sammensatt tilbud.
\end{itemize}

\subsection{Eksamensrelevans}
\begin{itemize}
  \item Definere digital innovasjon (Yoo et al.\ 2010) og forklare
    rekombinasjon.
  \item Forklare de tre egenskapene ved digital teknologi og hvorfor de
    muliggjør innovasjon.
  \item Forklare Christensens disrupsjonsteori.
  \item Vurdere et konkret initiativ som sustaining vs.\ disruptiv.
  \item Knytte til \autoref{sec:lec04} (DT vs.\ ITEOT) -- en disruptiv
    innovasjon medfører ofte digital transformasjon; en sustaining
    innovasjon er ofte IT-enabled.
\end{itemize}
